\section{总结与延伸}

\subsection{研究范式的转变}

杨植麟指出，当前 AI 研究范式与 10 年前截然不同。过去侧重"发表新 idea"，但缺乏严格实验的支撑难以得出可靠结论。如今，有了 scaling ladder 和丰富的 benchmark 体系，可以在不同规模上验证想法，从而对"古老"的技术（如优化器、注意力机制、残差连接）做出自信且扎实的改进。

\subsection{三个维度的乘法效应}

\begin{importantbox}{从加法到乘法}
三个 scaling 维度不是简单叠加，而是相互增强的乘法关系：
\begin{itemize}
    \item Adam (2014) $\rightarrow$ Muon-Clip：更好的 token efficiency
    \item Full Attention (2017) $\rightarrow$ Kimi Linear：更好的长上下文
    \item Residual Connection $\rightarrow$ Attention Residue：更高效的深度 scaling
\end{itemize}
将这些改进相乘，可以获得远超单项改进之和的整体提升。
\end{importantbox}

\subsection{开源社区的未来}

杨植麟对开源社区的未来持乐观态度，认为在架构和优化方面将持续涌现突破性改进。Agent Swarms 不是终点——scaling 的新维度会不断出现，推动开源模型与闭源模型的持续竞争。

\subsection{拓展阅读}

\begin{itemize}
    \item Kaplan et al., "Scaling Laws for Neural Language Models" (2020) --- 经典 scaling law
    \item Muon Optimizer 论文 --- Kimi 团队的二阶优化器工作
    \item Kimi Linear 技术报告 --- Delta Attention 的详细公式推导
    \item Attention Residue 技术报告 --- 深度维度注意力的完整实验
    \item He et al., "Deep Residual Learning for Image Recognition" (2015) --- ResNet 原始论文
\end{itemize}

%% --- Fin del contenido --- %%

\end{document}
